\scanpage{145}
% lines 1–8
Notons que l'image de $\unsure{\eta|y}$ dans $H^1(S,J)$, par fonctorialité, n'est autre que la classe \deleted{de} $\underline{\mathrm{Pic}}^1_{X/S}$ (considéré comme torseur sous $J$). Supposons cette classe nulle (ce qui est le cas si $X/S$ admet une section). \deleted{\unsure{p. ex.}} Alors la suite exacte
\[
H^0(S,J)\longrightarrow H^1(S,T)\to H^1(S,E)\unsure{-|\to}H^1(S,J)
\]
montre que \deleted{\unsure{le}} $\eta$ est dans $\operatorname{Im}H^1(S,T)$.

% lines 9–13
Notons qu'on a d'ailleurs une suite exacte \inserted{(si $Y-S$ surjectif)}
\[
0\to G_{m_S}\to\prod_{Y/S}G_m\to T-0 \tag*{($*$)}
\]
qui donne une suite exacte \deleted{\unsure{écrite}}
\[
\unsure{0|\cdot}\to H^0(S,\underline O_S)^{\unsure{*}}\to H^0(Y,\underline O_Y)^{*}\to H^0(S,T)-\mathrm{Pic}(S)\to\mathrm{Pic}(Y)-H^1(S,T)\to\mathrm{Br}'(S)\to\mathrm{Br}'(Y),
\]

% lines 14–21
Lorsque la suite exacte ($*$) splitte, p. ex. lorsque $Y$ admet une section sur $S$, alors on trouve
\[
\deleted{H^1(S,T)=}\qquad H^1(S,T)\simeq\mathrm{Pic}(Y)/\operatorname{Im}\mathrm{Pic}(S)
\]
et lorsque $Y=\coprod_{i\in I} g_i(S)$ comme dessus, alors \deleted{\unsure{on peut}} $\mathrm{Pic}(Y)\simeq\mathrm{Pic}(S)^I$ et on trouve
\[
H^1(S,T)\simeq\mathrm{Pic}(S)^I/\underset{\longrightarrow\,\text{par diagonale}}{\operatorname{Im}}\mathrm{Pic}(\deleted{\unsure{X}}S).
\]
Si p. ex. $\mathrm{Pic}(S)=0$, on trouve $\eta=0$, à fortiori les $\xi_n$, $\xi_l$ sont nuls.
% archive sheet number: 144

